\section{System Architecture}
\label{sec:architecture}

\subsection{Component view}

The platform is a small set of cooperating processes (Figure~\ref{fig:arch-component}).
The browser hosts the React SPA, which talks HTTPS to the Vercel CDN for
static content and to the Express API on Render for everything dynamic. The
API is the only process that holds Supabase service-role credentials; the
client never sees Supabase directly. A second outbound HTTP path connects
the API to the Anthropic API for LLM calls.

\fig{arch-component}{0.92}{Component view of the Spotlight platform.}

The choice of \emph{which} dynamic state lives where is deliberate:

\begin{itemize}[leftmargin=*]
  \item \textbf{Durable, multi-tenant state} (projects, votes, grades,
        audit log, AI insights cache, vote flags) lives in Supabase
        Postgres. It survives process restarts and is the single point of
        truth.
  \item \textbf{Hot, derivable state} (cached query results) lives in an
        in-process LRU. Wrong answers from this cache are bounded by the
        TTL: 60\,s for the project list, 30\,s for winners, 5\,s for vote
        counts.
  \item \textbf{Ephemeral fan-out state} (the set of connected SSE clients)
        lives in a process-local \texttt{Set}. A multi-instance deployment
        would need Redis pub/sub here, but a single Render dyno is enough
        for the event-day load.
  \item \textbf{Anomaly-detector ring} (last 5 minutes of votes) is also
        process-local. Persistence is intentionally absent: a flagged vote
        is durable in \texttt{vote\_flags}; the ring only feeds the scoring
        function.
\end{itemize}

\subsection{Deployment view}

Figure~\ref{fig:arch-deployment} shows the deployment target. The
frontend bundle is uploaded to Vercel on every push to \texttt{master}; the
Render service rebuilds on the same trigger. Supabase is provisioned
manually once.

\fig{arch-deployment}{0.78}{Deployment view.}

\subsection{Tech stack rationale; decision records}

Each major technology choice is captured below in the form of a lightweight
architecture decision record (ADR). The format is: \emph{decision}, \emph{context},
\emph{consequences}.

\begin{ncbox}{ADR-1: React 19 SPA + HashRouter}
\textbf{Context.} Vercel's free tier is happiest with a single static
\texttt{index.html}; we have no SEO requirement; deep linking inside
admin/judge surfaces is needed.\\
\textbf{Decision.} Ship the frontend as a Vite-built SPA and route with
\texttt{HashRouter}. Avoid SSR.\\
\textbf{Consequences.} Cold-load is the price of one bundle, but every
subsequent navigation is local. We avoid the operational complexity of
Next.js / SSR for a system whose first paint is rarely the bottleneck.
\end{ncbox}

\begin{ncbox}{ADR-2: Express 4 + Supabase, no ORM}
\textbf{Context.} An ORM (Prisma, Drizzle) would give us migrations and
typed queries, but it adds a build step, a generation step, and at the
scale of $\approx10$ tables it is overkill.\\
\textbf{Decision.} Use the Supabase JS client directly for parameterised
queries; keep migrations as plain \texttt{.sql} files in
\texttt{server/migrations/}.\\
\textbf{Consequences.} Loss of compile-time query checking, mitigated by
zod runtime validation at the controller boundary. Migrations remain
auditable as plain SQL.
\end{ncbox}

\begin{ncbox}{ADR-3: SSE over WebSocket}
\textbf{Context.} Real-time push from server to client is one-way and
infrequent (a few events per minute, peaking at $\sim$1 per second during
busy stretches). WebSockets need their own protocol upgrade and are
sometimes blocked by corporate proxies.\\
\textbf{Decision.} Use Server-Sent Events with a watchdog reconnect on the
client and a polling fallback if the env flag \texttt{VITE\_USE\_SSE} is
false.\\
\textbf{Consequences.} Native browser auto-reconnect, simpler server (just
\texttt{Content-Type: text/event-stream}), HTTP/1.1-only on Render which
costs us no concurrency at our scale.
\end{ncbox}

\begin{ncbox}{ADR-4: In-process LRU instead of Redis}
\textbf{Context.} The traditional answer is Redis for any shared cache.
But we run a single backend instance; an external cache adds a network
hop and another deployment artefact for negligible gain.\\
\textbf{Decision.} \texttt{lru-cache} in the Express process. Wrap reads in
\texttt{cache.getOrSet(key, ttl, loader)}. Invalidate on the writes that
matter.\\
\textbf{Consequences.} Cache is per-process, so a second Render instance
would have its own. Fine for the current footprint; if we scale out, swap
the implementation behind \texttt{cache.getOrSet} for a Redis-backed one.
\end{ncbox}

\begin{ncbox}{ADR-5: HMAC vote token, not a CAPTCHA}
\textbf{Context.} Anonymous voting needs some \emph{proof of presence} to
defend against curl loops, but we do not want CAPTCHA UX in a 90-minute
live show.\\
\textbf{Decision.} Two-step vote: \texttt{POST /votes/init} returns a
short-lived HMAC token bound to the browser token; \texttt{POST /votes}
requires the token. Combined with rate limits and the anomaly detector.\\
\textbf{Consequences.} A naive script that only POSTs to \texttt{/votes}
fails. A determined attacker who scripts both calls is still subject to
rate limits and detector exclusion. Acknowledged residual risk: a real
botnet with rotating IPs and headless browsers; mitigated by exclusion
from People's Choice rather than refusal of the vote.
\end{ncbox}

\begin{ncbox}{ADR-6: Anthropic Claude Haiku 4.5 with prompt caching}
\textbf{Context.} The advisory insight panel needs an LLM. Latency
matters (judges look at it during the show); cost should not balloon if
the panel is regenerated.\\
\textbf{Decision.} Use \texttt{claude-haiku-4-5-20251001} for speed; declare
the rubric system prompt with \texttt{cache\_control: \{type: 'ephemeral'\}};
cache the parsed JSON in \texttt{ai\_insights} keyed on \texttt{project\_id}.\\
\textbf{Consequences.} Per-call cost approaches the prompt-caching read
rate after the first call. A failed call (rate limit, network) falls back
to a deterministic stub so the UI never blanks.
\end{ncbox}

\begin{ncbox}{ADR-7: Pino structured logs, no APM}
\textbf{Context.} Datadog / New Relic / Sentry are overkill for a single
event. We still need to debug in production.\\
\textbf{Decision.} Use \texttt{pino} with \texttt{pino-http}; every request
gets a UUID via \texttt{X-Request-ID} or generated on the server, attached to
every log line and returned in error responses.\\
\textbf{Consequences.} Cheap, durable, greppable. We can pipe to a hosted
log collector later by adding a transport line.
\end{ncbox}

\subsection{Cross-cutting concerns}

\paragraph{Authentication and authorization.} A judge logs in, receives a
JWT signed with \texttt{JWT\_SECRET}, and stores it in localStorage. Every
protected request sends \texttt{Authorization: Bearer <jwt>}; the
\texttt{protect} middleware verifies the token, fetches the judge row,
and attaches \texttt{req.judge} including \texttt{isAdmin}. The
\texttt{adminOnly} middleware short-circuits non-admin requests with 403.

\paragraph{Validation.} Every controller route declares a zod schema in
\texttt{server/middleware/validate.js}. Invalid payloads are rejected with
400 before any DB call.

\paragraph{Error handling.} \texttt{asyncHandler} wraps every async
controller. Errors propagate to a global \texttt{errorHandler} that logs
with the request id and emits a clean JSON response.

\paragraph{Caching.} The \texttt{cache.getOrSet} helper is the only entry
point. Writes that change a cached view (\texttt{castVote},
\texttt{submitGrade}, \texttt{updateProject}) explicitly invalidate the
relevant key prefixes.

\paragraph{Real-time fan-out.} \texttt{eventBus.broadcast(event, payload)}
is the only emitter. Subscribers connect via \texttt{GET /api/session/stream}
and the bus heartbeats every 25 s to keep proxies from idling the socket.
